\section{总结与结论}

本文扩展了文献 \cite{Orginos:2017kos} 对部分子伪分布及约化伪 ITD 格点数据所作的领头对数分析。为此，我们纳入了约化伪 ITD 单圈层次的近期结果 \cite{Radyushkin:2017lvu}（另见 \cite{Ji:2017rah,Izubuchi:2018srq}）。

我们发现该修正包含一个大项，导致与 LLA 相比的重要数值变化。大修正之所以出现，是因为最重要图形所涉及的有效距离远小于名义距离 $z_3$。这改变了重标度关系 $\mu = k/z_3$ 中的系数 $k$（在我们的特定 ITD 情形从 $k_{\rm LLA} \approx 1$ 变到 $k \approx 4$），该关系允许把伪 PDF ${\cal P} (x,z_3^2)$ （近似地）转换为 $\overline{\rm MS}$ PDF $f(x,\mu^2)$。

重标度关系可作为快速估计和半定量分析的有益指引，而 $\overline{\rm MS}$ ITD 也可以应用精确单圈公式直接构建。利用它，我们用 \mbox{$0\leq z_3 \leq 4a$} 区域的数据得到了 \mbox{$\mu=1/a \approx 2.15$ GeV} $\overline{\rm MS}$ 标度处的 ITD ${\cal I} (\nu, \mu^2)$。

我们发现此标度下的 ${\cal I} (\nu, \mu^2)$ 接近 $z_3\sim 4a$ 的约化伪 ITD ${\mathfrak M} (\nu,z_3^2)$。由于 $a\leq z_3 \leq 4a$ 区域的所有数据与 $z_3=4a$ 处相差不大（见图~\ref{M16}），把 ${\mathfrak M} (\nu,z_3^2)$ 数据转换为 ${\cal I} (\nu, 1/a^2)$ 不涉及大的变化，也就是说，用约化伪 ITD ${\mathfrak M} (\nu,z_3^2)$ 表示 $\overline{\rm MS}$ ITD ${\cal I} (\nu, \mu^2)$ 的微扰展开是可控的。形式上的原因是，此情形下的大修正可以被吸收进依赖 $z_3^2$ 的演化项中，剩余修正很小。

唯象上，如此提取的 PDF 在大 $x>0.1$ 区域遵循全局拟合所给 PDF 的走势，但不重现其在 $x<0.1$ 区域的奇异行为。后者通常与 $\rho$ 介子雷杰轨迹的 $x^{-0.5}$ 图样相联系。既然 $\rho$ 介子本质上是 $\pi \pi$ 系统中相当窄的共振，在 $\pi$ 介子重达 600 MeV 的格点模拟中不应期望准确重现 $\rho$ 介子的性质。因此可以希望，采用物理 $\pi$ 介子质量的模拟会在小 $x$ 区域给出与全局拟合更好的符合。这一期望得到近期在物理 $\pi$ 介子质量下用准 PDF 格点模拟提取 $q_v (x)$ 的工作 \cite{Alexandrou:2018pbm,Chen:2018xof} 的支持。
